\section{Conclusion}

We built a working urban flood twin for 16 Indian city areas out of free satellite data, on free
computing, and put it on the public internet. That part is engineering. The part we think matters
more is what happened when we measured it properly.

A flood overlap score on this task is uninterpretable on its own. A baseline with no physics, no
rainfall and one fixed parameter scores 0.29 to 0.67 across 15 areas, beating every physics-model
score we can find published, and beating the agreement between two radar passes of the same city in
all fifteen. Any CSI in this literature, ours included, should be read next to those numbers or not
read at all.

Once that scale is fixed, the rest of the results change shape. Our physics twin scores 0.054 and a
learned model on the identical rasters scores 0.288, which refutes our own earlier claim that
terrain fundamentally lacks the information. But that same learned model ties the parameter-free
baseline exactly, so its value is not accuracy. Its value is that it runs on cities with no radar
history at all, and held out one whole region at a time that is worth 0.147, rising to 0.175 with a
threshold correction that costs nothing and never looks at a test label. The one place it nearly
fails is a coast it has not seen, and we showed that is a coverage problem rather than a hard
limit, because adding a single coastal city lifted held-out coastal skill by 84\,\%.

Three things that ought to have helped did not. Model size across a 35-fold range did nothing.
Grid resolution from 60 to 30\,m did nothing. Rainfall from free reanalysis did nothing at best and
actively hurt at worst, shown four independent ways with two alternative explanations refuted.

Meanwhile planning works despite the location uncertainty, because we measure every intervention by
re-simulation instead of assuming it. A dense street-routed drain network cuts simulated street
flooding by 80.9\,\% in Patna and costs less per cubic metre than the sparse alternative. A metered
recharge planner banks 6.3 to 16.5\,\% of a storm's runoff in six depleted Karnataka towns against
0.7 to 3.3\,\% on the flood plain, and cuts flooding at the same time. An inferred drainage field
puts a falsifiable ceiling on Mumbai's whole-city drainage near 5.4 million cubic metres per storm,
which sits on the correct side of the 4.7 million cubic metres per day the city measured itself
pumping. And a graph network over the streets makes all of it fast enough to route a person around
the water in four milliseconds.

We have reported the failures beside the successes throughout, including three corrections against
our own published claims, a negative result against 83 crowdsourced depth measurements, a drainage
inversion that improved the quantity it was trained on and did nothing for the one that matters,
and a learning loop whose gate has never fired because nobody has filed a report yet. We would
rather the record be useful than flattering.

\section*{Reproducibility}
Code, all per-area artifact bundles including the Sentinel-1 water masks for the 15 areas that have
them, road graphs, replay buffers and ranked site lists, together with 253 offline tests and
step-by-step runbooks, are public at
\url{https://github.com/Maverick-Ansh/project-varuna}. The live dashboard is at
\url{https://project-varuna-iota.vercel.app} and its API at
\url{https://anshvivek-varuna-floodtwin.hf.space}.

A new city builds with one script on a free GPU and serves from the committed bundle. Every headline
number in this paper is regenerable from a committed artifact. The learned-baseline dataset is
rebuilt by three builders, each with a verify mode that reports agreement with the committed copy
rather than asserting it. The headline is unchanged on a fully rebuilt dataset, at
$0.288 \pm 0.008$ against $0.288 \pm 0.007$. The trivial-baseline builder reproduces all previously
published values before extending them, which is how the fifteenth area was added to
Table~\ref{tab:trivial} without disturbing the other fourteen. Sentinel-1 masks for a new area are
fetched by a script that needs one interactive Earth Engine login and ranks candidate dates by
antecedent rainfall, so scenes are chosen for having weather in them rather than by hand.

\section*{Acknowledgment}
This work uses FABDEM, ESA WorldCover, JRC Global Surface Water, Copernicus Sentinel-1,
OpenStreetMap, Open-Meteo, NASA Black Marble, SoilGrids, HydroSHEDS and Central Ground Water Board
data, all of which are free to use. The crowdsourced depth observations come from the public
mumbaiflood.in service. The system runs on free tiers provided by Google Colab, Kaggle, Hugging Face
and Vercel. None of this work would be possible if those datasets and services were not open.
